\section{Kouzlení a čarování}

\subsection{Pomůcky}

Kouzelnická hůlka je dřevěná, kulatá 30 cm dlouhá tyčka. Celá 
černě lakovaná, jen 6 cm obou konců je bílých.

\subsection{Kouzla humorná}

\nadpis{Fakir}{cokouzla}
Palec překryji šátkem a začnu do něj bodat jehlice a 
špendlíky. Pak je zase vytahám, palec odkryji a vše je v~naprostém 
pořádku.

Palec z~rukavice něčím vycpaný je jemně nastehován 
na šátek. Když palec přikrývám, pod kapesníkem palec ohnu a do 
prstů sevřu spodní konec látkové napodobeniny. Napodobenina musí 
být vždy schovaná v~šátku, divák ji nesmí vidět.


\nadpis{Jasnovidec}{cokouzla}
Řekneme, že uhádneme, co si kdo myslí. Do řady postavíme 
tři dobrovolníky a připravíme tři mince. Prvnímu dobrovolníkovi 
řekneme: ,,Dám vám teď minci, pevně ji sevřete v dlani.`` A dáme 
mu minci. U druhého postup zopakujte. Třetímu však nic neříkejte, 
jen viditelně držte poslední minci. Divák otevře dlaň a nastaví 
ruku. Ukažte na něj hůlkou a řekněte: ,,Tento divák si myslí, 
že také dostane minci.``


\nadpis{Jasnovidka}{cokouzla}
Kouzelnice požádá několik diváků, aby napsali na papír 
několik slov. Oni pak odevzdají asistentce, která je složí, vloží 
do klobouku a odevzdá jasnovidce. Ta papíry vytahuje, čte slova 
a říká, zda je psala žena či muž.

Trik spočívá ve skládání papíru. Papír lze nejprve přehnout 
ve svislé nebo ve vodorovné ose. Podle toho jasnovidka pozná 
pisatele.


\nadpis{Jasnovidec}{cokouzla}
Diváci se domluví na nějakém slově a já je uhádnu.

Kouzelník ukazuje náhodně na diváky a ti mají za úkol říkat 
různá slova. Kouzelník je však domluven s~pomocníkem, 
který říká slova, jejichž první písmeno je vždy obsaženo v hledaném 
slově. Je-li slovo celé, řekne pomocník předem dohodnuté slovo.

Př. Slovo=řemen, pomocník hlásí Řád, Eskymák, Mlha, Ešus, Noha.


\nadpis{Kuličky}{cokouzla}
Z papíru 3$\times$3 cm utvořte tři kuličky a položte si je do 
žlábku mezi dvěma prsty. Pak foukněte a všechny kuličky spadnou. 
Teď se zeptejte, která kulička má zůstat na svém místě při dalším 
fouknutí. Zamávejte nad ní hůlkou,přidržte ji a foukněte. Diváci
pochopí žert, vy však buďte ochoten pokus opakovat s jinou kuličkou.


\nadpis{Kolik je 7$\times$7}{cokouzla}
Utvoříme si kartičku 7$\times$7 a rozstříháme ji podle prvního
obrázku. 7$\times$7 je teď 49. Poskládáme-li však rozstříhanou
kartičku podle druhého obrázku, přesně jeden čtvereček bude
chybět. 7$\times$7 je tedy 48.\\
\includegraphics{kolik7x7}


\nadpis{Pravopisný chyták}{cokouzla}
Ať se přihlásí ten, kdo si myslí, že zná dobře český pravopis. 
Nadiktujeme mu pak k napsání tuto větu: Kočí zvedl bič a zvolal: 
Hyje, koníčky...  (na konci mlaskneme jazykem jako na koně). 
Mlasknutí se ještě nikomu nepodařilo napsat a ani nepodaří.


\nadpis{Silné nervy}{cokouzla}
,,Kdo z vás vydrží pod stolem do kterého 3x za sebou prudce 
udeřím dlaní?'' Určitě se nikdo najde a vleze si pod stůl. Ty 
pak udeř do stolu ale jen dvakrát. ,,Tu třetí ránu dám až za 
týden, jsem sám zvědav, jestli odvážlivec tam tak dlouho vydrží 
!''


\subsection{Kouzla vážná}

\nadpis{Prostupnost hmoty}{cokouzla}
Mám šňůrku a na jejích koncích korálky. Vezmu sklenici 
s~velkým uchem dnem vzhůru. Šňůrku položím na levou ruku 
tak, že její střed spočívá mezi nataženým palcem a ukazováčkem 
a tvoří mezi nimi proužek. Na šňůrku položím ucho sklenice, zatlačím, 
korálky se mírně zvednou. Pak sklenici pustím a ona zůstane viset.

Když tlačíme ucho proti šňůrce a tím vytahujeme korálky, 
vyprostíme palec, takže stuha visí je na ukazováku. Poté vsuneme 
ucho mezi palec a šňůrku a vsuneme palec uchem zase pod šňůrku. 
Tlačíme dál, až se nám dostanou korálky do dlaně. Teď sklenici 
pustíme. Jeden korálek se provleče uchem, šňůrka se napne a sklenice 
zůstane viset.


\nadpis{Čtenář myšlenek}{cokouzla}
Oznámíte všem, že jste čtenář myšlenek, o čemž je hned 
přesvědčíte. Vyzvete je, aby každý napsal na list papíru nějaké 
přísloví nebo libovolnou větu o pěti slovech (ne více), dal ji 
do obálky a zalepil. Potom obálky seberte, jednu po druhé přikládejte 
na čelo, soustředíte se, tváříte se důmyslně a k úžasu přítomných 
čtete věty v nich obsažené.

Je třeba abyste byli předem smluveni s některým z hráčů. 
Když přiložíte první obálku na čelo, nevíte sice co je v ní napsáno, 
ale znáte větu, kterou jste si předem smluvil se svým tichým 
společníkem a tuto větu řeknete jako první. Když pak obálku ,,pro 
kontrolu'' otvíráte, seznámíte se s obsahem věty, kterou 
napsal jiný účastník a přisoudíte ji další obálce. Po skončení 
celé atrakce se mohou účastníci přesvědčit, že všechny věty, 
které jste přečetl, jsou skutečně na lístcích napsány. Při otvírání 
obálek je nutno dát bedlivý pozor, abyste obálku se smluvenou 
větou otevřeli jako poslední. Je dobře, když obálky i lístky 
papíru jsou stejné, aby některý z diváků nepoznal trik. Pokud 
chceme ušetřit obálky, musí všichni hráči svůj papír složit stejným 
způsobem a to tak, aby nebyl vidět text.


\nadpis{Sirky ve vzduchu}{cokouzla}
Ukážu ruce, sáhnu do vzduchu a objevím sirku, hodím 
ji do klobouku a už vytahuji za zády další sirku, kterou opět 
házím do klobouku a to už... Pak z klobouku vysypu celou hrst 
sirek

Z průhledné lepicí pásky si zhotovím lepicí prstýnek a nasadím 
na prostředník. Jakmile ukážu ruce, nalepím na prstýnek ze strany 
nehtu sirku. Jsou-li prsty nataženy, sirka není vidět. Hmátnu 
tedy do vzduchu, skrčím prsty a objevím sirku. Objevenou sirku 
přidržím palcem. Ruku vstrčím do klobouku, kde prsty narovnám. 
Tím se sirka zase schová a já ji mohu zase někde najít. Nakonec 
sirku palcem odlepím, hodím tentokráte doopravdy do klobouku 
a z klobouku vysypu sirky, co jsem tam měl již nachystány.


\nadpis{Dámy utekly}{cokouzla}
Ukážeme několik karet a mezi nimi dvě dámy. Řekneme, 
že dámy nemají rády, když to a to. Pak karty složíme a převážeme 
gumičkou. Uděláme to a to. Poté gumičku odstraníme a vida, dámy 
jsou pryč. Pak je někde objevíme.

K tomuto kouzlu je nutné si připraviti preparovanou kartu. 
Z obou dam odstřihneme horní pětinu karty a slepíme je k sobě. 
Mezi ostatními kartami (ukazujeme sloupec karet) vypadají normálně. 
Když ale karty skládáme, přehneme preparované dámy a schováme 
je v dlani. Jakmile sáhneme do kapsy pro gumičku, necháme dámy 
v kapse. Před představením někde uschováme dvě normální dámy, 
které pak náhodou objevíme.


\nadpis{Kde je těžiště}{cokouzla}
Kouzelník vezme dlouhou krabičku a z malé části ji položí 
na sklenici, krabička však nespadne.

U jednoho konce krabičky je připevněno olůvko, které přesunuje 
těžiště, takže se sklenice nepřeváží.


\nadpis{Kdo postaví krabičku}{cokouzla}
Prázdná krabička od sirek je položena na kraji stolu 
tak, že ji kousek přečnívá desku stolu. Kouzelník ji zespodu 
jedním prstem postaví. Ukáže krabičku, že je zcela normální. 
Když krabičku zkouší postavit divák, vždy se mu převáží dozadu.

Krabička na stole leží vždy obrázkem nahoru. Je-li dno její 
zásuvky blíže k desce stolu, krabičku lze postavit, je-li však 
zásuvka zasunuta obráceně, krabička se vždy převáží.


\nadpis{Matematický omyl}{cokouzla}
Divák má sčítat po řadě čísla 1000 + 30 + 1000 + 30 + 
1000 + 30 + 1000 + 10. Téměř všichni odpoví, že výsledek je 5000, 
ačkoli je 4100.

Divák prostě podlehne pravidelnosti a svůj omyl si neuvědomí.


\nadpis{Mince je pryč}{cokouzla}
Položíme minci na táce, přikryjeme šátkem a obejdeme 
diváky. Každý vsune ruku pod šátek a řekne, zda tam mince skutečně 
je. Potom uděláme abrakadabra a zvedneme šátek. Mince zmizela. 
Obdobným postupem se mince na tácku opět objeví.

Ke kouzlu je zapotřebí pomocníka, jemuž jako poslednímu tácek 
ukazujeme. Ten minci v první fázi vezme, schová a řekne, že tam 
je. Ve druhé fázi minci na tácek položí a řekne (jako všichni 
ostatní), že tam mince není.


\nadpis{Najdu správnou kartu}{cokouzla}
Karty jsou seřazeny do 3 sloupců po 7 kartách. Divák 
si vybere kartu a ukáže na sloupec, ve kterém se nachází. Totéž 
pak udělá ještě 3x. Pak mu kouzelník slavnostně ukáže jeho kartu.

Karty je nutné sbírat po sloupcích tak, aby divákem označený 
sloupec byl vždy uprostřed. Po třetím posbírání je karta je jedenáctá 
ať se počítá z kterékoli strany.


\nadpis{Očarovaný řetěz}{cokouzla}
Dvojitý řetěz držím ve svislé poloze. Druhou rukou vezmu 
kroužek, pustím jej a on propadne a zastaví se dole jako poslední 
článek řetězu.

Asi ze třiceti kroužků sestavíme řetěz následovně. Do jednoho 
kroužku navlékneme dva další. Pak řetěz přetočíme o 90$^\circ$ 
doprava. Do zadního kroužku navlékneme další kroužek a oběma 
(zadním i předním) provlékneme druhý kroužek. Řetěz opět přetočíme 
o 90$^\circ$ doprava a pokračujeme v sestavování. Do zadního 
kroužku... Poslední kroužek je navléknut do obou kroužků předchozího 
článku.

Palec a ukazovák levé ruky pevně drží jednotlivý horní kroužek. 
Palec a ukazovák pravé ruky uchopí pravý kroužek pravého páru 
(je v něm navléknut pouze jeden kroužek druhého páru) Horní jednotlivý 
kroužek pustíme a ten zdánlivě prostoupí celým řetězem a umístí 
se jako poslední dole. Přitom řetěz přetočíme zase o 90$^\circ$ 
a horní kroužek uchopíme levou rukou. Palec a ukazovák pravé 
ruky uchopí pravý kroužek pravého...


\nadpis{Přimražená sklenice}{cokouzla}
Na sklenici úplně plnou vody přiložíme papír a zrcátko. 
Pak foukneme ,,studený vítr`` a posečkáme. Sklenice za určitou 
dobu přimrzne k zrcátku.

Papír je savý piják, voda převařená. Papír vysává vodu, tím 
vzniká podtlak a sklenice drží.


\nadpis{Sirka dělá stojku}{cokouzla}
Kouzelník má v ruce sirku, postaví ji hlavičkou na podložku, 
pronese kouzelné slovo a sirka stojí.

Potají si nasliníme nehet, ten pak v nestřeženém okamžiku 
třeme o hlavičku sirky, čímž ji navlhčíme. Takováto sirka vydrží 
stát na vhodné podložce (papír, látkový ubrus) velmi dlouho.


\nadpis{Sklenice propadne stolem}{cokouzla}
Kouzelník vezme sklenici, přikryje je ji šátkem. Pak 
nakreslí na stůl křížek a položí sklenici na něj. Sklenice však 
propadne stolem, spadne na zem a rozbije se.

Kartónové kolečko stejného průměru jako vrchol sklenice nastehujeme 
na šátek. Když bereme sklenici pod šátkem, chytneme sklenici 
i kartónek. Ve chvíli, kdy kreslíme křížek, pustíme sklenici 
do klína (díky šátku nepozorovaně). Kartónek stále držíme, takže 
to vypadá, jako by se nic nestalo. Pak už jen stačí, když pokládáme 
,,jako sklenici`` na stůl, pustit z klína sklenici na zem.


\nadpis{Voda mění barvu}{cokouzla}
Ve sklenici je červená voda. Sklenici přikryjeme šátkem, 
zakouzlíme. Když zvedneme šátek, je ve sklenici voda čirá. Přelijeme 
vodu z jedné sklenice do druhé a vida --- voda je zase červená.

Ve sklenici je vložka z červené hmoty vysoká přesně na výšku 
hladiny. Když snímáme šátek, vyjmeme zároveň i vložku, takže 
voda je bílá. Ve druhé sklenici je KMnO$_4$, ten vodu obarví opět 
na červeno.


\nadpis{Výroba zlata}{cokouzla}
Podivným činěním vytvoříme zlato.

Roztok octanu olovnatého v jedné baňce a jodidu draselného 
v baňce druhé zahříváme na teplotu varu. Slijeme obě tekutiny 
barvy čisté vody dohromady. Jakmile začne směs chladnout, objeví 
se na dně krystalky ,,zlata''.


\nadpis{Zmizelé mince}{cokouzla}
Na stole leží mince. Vezmeme plechový tácek, shrneme 
na něj zmiňované mince. Mince při dopadu na tácek zvoní, ale 
když jej ukážeme divákům, nic na něm není.

Asi 10 cm pod úrovní desky stolu je plech. Mince tedy zvoní, 
jako kdyby dopadaly na tácek, ve skutečnosti však dopadají na 
plech.


\nadpis{Modrý --- červený}{cokouzla}
Karetní hru rozevřu do vějíře a ukážu z~obou stran. 
Diváci vidí karetní hru s~modrým rubem. Pak karty shrnu 
do balíčku, modrou kartu odeberu, červenou přidám, karty zase 
rozevřu a už mají rub červený.

Karty preparujeme následujícím způsobem. Oddělíme ruby a 
úhlopříčně je rozstřihneme a pak zase slepíme. Lepíme však tak, 
že na jedné kartě je polovina červeného a polovina modrého rubu. 
Slepené karty zatížíme. Dvě karty takto nepreparujeme. Když karty 
rozevřeme, ukážeme vždy jednu polovinu rubu. Jen horní (jedna 
z~nepreparovaných) je vidět. Pak karty složíme a vodorovně 
odložíme. Poté sejmeme modrou nepreparovanou a nahradíme ji červenou, 
rovněž neupravenou kartou. Mezitím je nutné karty otočit. Když 
vějíř zase rozevřeme, vidí diváci tentokráte druhou polovinu 
rubu preparovaných karet.


\nadpis{Colombovo kouzlo}{cokouzla}
Vyzvěte diváka, aby si myslel číslo od 1 do 4 a pak vám 
je řekl. Na důkaz své jasnozřivosti vyzvěte diváka, aby se podíval 
na určité skryté místo, kde nalezne cedulku: ``Věděl jsem, že 
si myslíte číslo...''.


Na různých skrytých místech jsou umístěny cedulky pro všechna 
přípustná čísla. Podle čísla, které si myslel, vyšle diváka kouzelník 
na místo s~patřičnou cedulkou.


\nadpis{Jasnovidec}{cokouzla}
Vybraný divák si bude myslet číslo od 1 do 4 (lze i více). 
Kouzelník se chvíli soustředí a pak řekne: ,,Už vím, jaké 
číslo si myslíte. Řekněte mi je.'' Divák řekne své myšlené 
číslo. ,,A teď se podívejte pod talíř,'' řekne kouzelník. 
Divák pod talířem nalezne cedulku: ,,Věděl jsem, že si 
myslíte číslo...''

Pokud by divák řekl, že si myslel jiné číslo, poslal by ho 
kouzelník podívat se pod jiný předmět. Kouzelník má totiž připraveny 
pod různými předměty cedulky pro všechna možná čísla. Předmětů 
by mělo být mnohem více než čísel.


\nadpis{Kapsa plná karet}{cokouzla}
Divák vytáhne z~balíčku 3 karty a zapamatuje si 
je. Kouzelník si karty vloží do kapsy, potom dvě vytáhne a požádá 
diváka, aby řekl jednu z~karet, které si vytáhl. Kouzelník 
potom vytáhne z~kapsy třetí kartu --- a je to přesně ta, 
jakou určil divák.

Kouzelník si karty vloží do kapsy a přitom si zapamatuje 
jejich pořadí. Potom vytáhne z~kapsy dvě jiné karty, které 
už v~kapse byly. Jakmile řekne divák své přání, bez problémů 
vytáhne patřičnou kartu.


\nadpis{Létající mince}{cokouzla}
Před zraky diváků postav na prázdnou sklenici talířek. 
Přes všechno pak přehoď šátek. Nyní vezmi do levé ruky korunu 
a dej ji do druhé ruky a zároveň ji vyhoď do výšky. Vezmi kouzelnickou 
hůlku a začni vysvětlovat, že mince lítá ve vzduchu a že jí hůlkou 
stáhneš přes talířek do sklenice. Hůlkou se dotkneš a uvolněná 
mince spadne s cinknutím na dno sklenice.

Při pokládání talířku na sklenici přidrž nenápadně na dně 
talířku přidrž korunu a přimáčkni ji mezi talířek a sklenici. 
Do levé ruky vezmi jinou korunu, naoko ji dej do druhé ruky a 
naznač její vyhození do vzduchu. Mezitím ukliď z levé ruky korunu 
do kapsy. Hůlkou se potom dotkni talířku na okraji, tím ho na 
druhé straně nadzvedneš a mince spadne do sklenice.


\nadpis{Matematický trik}{cokouzla}
Kouzelník předloží kartičku s~čísly. Divák si vybere 
jedno číslo a zatrhne všechna čísla v řádku a sloupci, kde se 
zvolené číslo vyskytuje (10). Pak si vybere číslo ještě nezatržené 
a postup opakuje (4). Součet zbylých čísel je vždy 40.


\nadpis{Pravdomluvné zrcadlo}{cokouzla}
Vyhlídni si předem svou oběť, na kterou napiš křídou 
na zrcadlo nějaký vtip. Vše pak smaž hadrem, takže sklo se jeví 
jako čisté. Pak přede všemi vyzvi svou oběť aby se podívala do 
zrcadla a řekla co vidí. Samozřejmě, že uvidí jen sama sebe. 
Řekni, že zrcadlo je kouzelné, ale ten co se do něj dívá se musí 
napřed otočit 3x doprava, pak 3x doleva a nakonec na zrcadlo 
3x dýchnout. Pak mu zrcadlo něco poví.


\nadpis{Rentgenový zrak}{cokouzla}
Kouzelník se otočí, diváci položí na stůl minci a přikryjí 
ji neprůhledným hrnečkem. Kouzelník se zamyslí a pozná, jaká 
to je mince.

Pomocník a zasvěcenec, který minci přikrývá, nasměruje ouško 
hrnečku směrem podle toho, jaká mince pod hrnečkem je.


\nadpis{Šáteček a svíčka}{cokouzla}
Na stolek připrav větší nezapálenou svíčku a pootevřenou 
krabičku sirek. Svíčku zapálíme a krabičku sirek zavřeme. Hořící 
svíčku chytíme jednou rukou a druhou vytáhneme z~plamene 
šáteček.

Šáteček je nasoukán v~krabičce sirek. Svíčku zapálíme 
a krabičku sirek zavřeme, tím se dostane šátek nepozorovaně do 
dlaně. Není pak už nic snazšího, svíčku chytit rukou se šátečkem 
a druhou jakoby s plamene vytáhnout šáteček.


\nadpis{Zlomená sirka}{cokouzla}
Divák zlomí sirku, kouzelník ji zabalí do kapesníku, 
zakouzlí a divákova sirka je opět celá a neporušená.

Předem (mimo zraky diváků) zasuneme do lemu kapesníku jednu 
sirku. Pak požádáme někoho z diváků, aby nám donesl zdravou zápalku. 
Tu pak před všemi dáme doprostřed kapesníku a poznačíme fixem, 
aby nemohlo dojít k záměně. Pak kapesník překládáme ovšem tak 
šikovně, abychom dali divákovy přes kapesník do ruky sirku v 
lemu. Požádáme ho aby sirku zlomil. Po té kapesník rozbalíme 
a vyndáme a divákům ukážeme opět zdravou sirku.


\nadpis{Živá kouzelnická hůlka}{cokouzla}
Po krátké přednášce o důležitosti kouzelnické hůlky, 
vezmi láhev a hůlku do ní strč. Po magických slovech a pohybech 
volné ruky sama vylézá z~láhve.

Hůlku však předem za jeden konec přivážeme za nit. Ten pak 
strkáme do láhve. Druhý konec nitě zavážeme k opasku kalhot. 
Tím, že pohybujeme lahví, hůlka vystupuje či zajíždí do láhve. 
